\section{Introduction}\label{sec:introduction}

A Hilbert function records the dimensions of a graded quotient; its graded
Betti numbers also record the relations among its generators and all higher
syzygies. The lex-plus-powers conjecture predicts one extremal ideal for
both kinds of data when the original ideal contains a regular sequence.
We prove its full Artinian form over \(\C\), and deduce the full
characteristic-zero statement for regular sequences of arbitrary positive
length by the established Artinian reduction and coefficient-field descent.

Let
\[
 S=\C[x_1,\ldots,x_n],\qquad
 2\le a_1\le\cdots\le a_n,\qquad
 P=(x_1^{a_1},\ldots,x_n^{a_n}),
\]
where each variable has degree one. Lexicographic order always has
\(x_1>\cdots>x_n\), and a lex segment starts with the \emph{greatest}
monomials. Suppose that a homogeneous ideal \(I\subseteq S\) contains a
homogeneous regular sequence \(f_1,\ldots,f_n\) with \(\deg f_i=a_i\).
Here regularity means that multiplication by \(f_i\) is injective on
\(S/(f_1,\ldots,f_{i-1})\) for each \(i\). Its full length makes
\(S/I\) Artinian: the complete-intersection Hilbert series is
\(\prod_i(1+z+\cdots+z^{a_i-1})\).
For each \(d\ge0\), prescribe
\begin{equation}
 J_d=L_d+P_d,\qquad \dim_\C J_d=\dim_\C I_d,
 \label{eq:prescription}
\end{equation}
where \(L_d\) is the largest lex-segment subspace with this property.
Part of the assertion below is that these dimensions can be attained
and that the prescribed spaces form an ideal.

For a standard graded polynomial ring \(A\) over a field \(F\), write
\[
 \beta^A_{p,j}(B)=\dim_F\Tor^A_p(B,F)_j,
 \qquad F=A/A_{>0}.
\]
The index \(p\) is the homological degree and \(j\) is the internal degree.

\begin{theorem}\label{thm:main}
For every \(n\ge1\), every ordered degree sequence
\(2\le a_1\le\cdots\le a_n\), and every homogeneous ideal
\(I\subseteq\C[x_1,\ldots,x_n]\) containing a regular sequence of these
degrees, the spaces in \eqref{eq:prescription} exist and
\(J=\bigoplus_{d\ge0}J_d\) is a homogeneous ideal. For this single ideal,
\[
 \beta^S_{p,j}(S/I)\le\beta^S_{p,j}(S/J)
 \qquad\text{for every \(p\ge0\) and \(j\in\Z\)}.
\]
There is no restriction on the degrees or number of the additional
generators of \(I\).
\end{theorem}

Thus Theorem~\ref{thm:main} resolves the Artinian lex-plus-powers
conjecture positively over \(\C\).
Its Hilbert-function assertion is the Artinian Eisenbud--Green--Harris
statement; the Betti inequalities are stronger.

The Artinian restriction can be removed from these two conclusions.
Caviglia and Maclagan proved the Artinian reduction for the
Eisenbud--Green--Harris conjecture \cite[Proposition~10]{CM08};
Caviglia and Kummini proved the corresponding reduction for
lex-plus-powers Betti numbers \cite[Section~4, Theorem~4.1]{CK14}.
Together with a finite coefficient-field descent, these reductions give
the following form over every characteristic-zero field.

\begin{corollary}[Eisenbud--Green--Harris and lex-plus-powers in characteristic zero]
\label{cor:characteristic-zero}
Let \(F\) be a field of characteristic zero, let
\[
 S_F=F[x_1,\ldots,x_n],\qquad
 1\le c\le n,\qquad 2\le a_1\le\cdots\le a_c,
 \qquad P_F=(x_1^{a_1},\ldots,x_c^{a_c}),
\]
with the standard grading and lexicographic order \(x_1>\cdots>x_n\).
Suppose that a homogeneous ideal \(I\subseteq S_F\) contains a homogeneous
regular sequence \(f_1,\ldots,f_c\) with \(\deg f_i=a_i\).
For \(d\ge0\), put
\[
 q_d=\dim_F I_d-\dim_F(P_F)_d.
\]
This integer lies between zero and the number of degree-\(d\) monomials
outside \(P_F\). Let \(J_d\) be the sum of \((P_F)_d\) and the span of the
\(q_d\) lexicographically greatest such monomials. Then
\(J=\bigoplus_{d\ge0}J_d\) is a homogeneous ideal, and for this single ideal
\[
 \dim_F(S_F/I)_d=\dim_F(S_F/J)_d\quad(d\ge0),\qquad
 \beta^{S_F}_{p,j}(S_F/I)\le\beta^{S_F}_{p,j}(S_F/J)
 \quad(p\ge0,\ j\in\Z).
\]
There is no restriction on the degrees or number of the additional
generators of \(I\).
\end{corollary}

Corollary~\ref{cor:characteristic-zero} includes non-Artinian quotients
and fields that are not algebraically closed. The case with no prescribed
regular sequence is the classical lex-ideal theorem cited below.
The division-coefficient construction itself remains the full-length
construction over \(\C\); Section~\ref{sec:characteristic-zero} proves
the corollary from Theorem~\ref{thm:main}.

\subsection{Context and previous results}

The Eisenbud--Green--Harris (EGH) Hilbert-function conjecture arose in
their work on higher Castelnuovo theory \cite[Section~4]{EGH}.
The lex-plus-powers conjecture, generally attributed to Charalambous
and Evans, strengthens that prediction to graded Betti numbers.
Francisco and Richert discuss its origins and identify its first
printed appearance in the work of Evans and Richert
\cite[Section~5]{FR}; see also the public formulations in
\cite[Section~1]{Richert} and
\cite[Conjectures~1.1 and 1.2]{CS}.
Without a prescribed regular sequence, the extremality of lex ideals
at fixed Hilbert function is the theorem of Bigatti, Hulett, and
Pardue \cite{Bigatti,Hulett,Pardue}.
The additional powers encode the degrees of the regular sequence and
give the sharper comparison studied here.
For ideals already containing these powers, Clements--Lindstr\"om
compression supplies the Hilbert-function comparison \cite{CL}.

The prescribed-powers Betti comparison developed through several stages.
Mermin, Peeva, and Stillman proved the comparison for ideals containing
the squares of the variables in characteristic zero
\cite[Theorem~1.1]{MPS}; Murai established it for strongly stable ideals
plus the prescribed powers over any field \cite[Theorem~1.1]{Murai}.
Mermin and Murai then proved the lex-plus-powers comparison for ideals
containing a monomial regular sequence, again over every field
\cite[Theorems~3.1 and 8.1]{MM}. They also isolated the reduction to a
Hilbert-matching monomial ideal with no smaller graded Betti numbers
\cite[Conjecture~10.2]{MM}. It is this reduction that we obtain from
division coefficients. Our proof of the subsequent monomial comparison
is included for completeness; its conclusion is not new.

For nonmonomial regular sequences, the Hilbert-function and Betti problems
have had different ranges of known cases. Abedelfatah proved EGH for full
regular sequences whose forms split into linear factors
\cite[Corollary~4.3]{Abedelfatah15}. Caviglia and Maclagan established
EGH for fast-growing degrees \cite{CM08}, and Caviglia and De Stefani
weakened the required growth to
\(a_i\ge\sum_{j<i}(a_j-1)\) for \(i\ge3\)
\cite[Theorem~A]{CDeS20}. These Hilbert-function results hold over any
field. The characteristic-zero Betti theorem of Caviglia and Sammartano
requires
\(a_i\ge1+\sum_{j<i}(a_j-1)\) for \(i\ge3\)
\cite[Theorem~3.8]{CS}. No such growth condition is imposed here.

More recently, Kuzmanovski \cite[Theorem~1.9]{Kuzmanovski} obtained
asymptotic Betti estimates in characteristic zero in a bounded degree range,
for fixed regular-sequence degree offsets and generator counts and
sufficiently large initial degree. That comparison uses the
reverse-lexicographic generic initial ideal of a lex-plus-powers ideal,
rather than the lex-plus-powers ideal itself as in Theorem~\ref{thm:main}.

Our treatment of the monomial comparison retains the
resolution-level mechanism of distractions
\cite[Theorem~2.19 and Corollary~2.20]{BCR} and the endpoint
decomposition for stable-plus-powers ideals
\cite[Lemma~3.7]{CS}, while organizing the last comparison through ranks
in the Koszul complex.

\subsection{The reduction that carries the Betti numbers}

The essential step is to find a monomial ideal \(M\supseteq P\) such that
\begin{equation}\label{eq:overview-bound}
 \Hilb(S/I)=\Hilb(S/M),\qquad
 \beta^S_{p,j}(S/I)\le\beta^S_{p,j}(S/M)
 \quad\text{for all \(p,j\)}.
\end{equation}
Matching Hilbert functions alone cannot give the second conclusion.

Theorem~1.1 of the companion article \cite{ArtinianEGH} supplies a field
\(k/\C\), a finite-dimensional central division \(k\)-algebra \(\Delta\), and commuting linear
forms \(t_1,\ldots,t_n\) in \(N=\Delta[x_1,\ldots,x_n]\). They satisfy
triangular power identities adapted to the \(f_i\), and their ordered
monomials form a left-\(\Delta\) basis of \(N\).
Section~\ref{sec:forms} states exactly the input needed here, including
the basis assertion for the full algebra rather than only its
complete-intersection quotient.

Let \(T=k[y_1,\ldots,y_n]\) act on \(N\) on the right by \(y_i\mapsto t_i\).
Then \(N\) is simultaneously left-\(S\) free and right-\(T\) free, the
latter of rank \(m=\dim_k\Delta\), with generators in degree zero.
A minimal \(S\)-resolution of \(S/I\) therefore gives a
\emph{minimal} \(T\)-resolution of \(N/IN\), multiplying every Betti number
by the same factor \(m\). With a monomial order adapted to the triangular
identities, the leading coefficient spaces of \(IN\) in the
\(t^\alpha\)-expansion are left ideals of \(\Delta\), hence either zero
or all of \(\Delta\). Their nonzero positions define \(M\); the
triangular identities ensure that \(M\) contains \(P\). A filtered
Koszul comparison then gives \eqref{eq:overview-bound} after cancelling
\(m\).
This resolution transfer is the additional ingredient beyond the
Hilbert-function reduction in \cite{ArtinianEGH}.

Starting from \(M\), we work with the monomials in the exponent box
\(0\le u_i<a_i\). A transfer replaces a factor \(x_\ell\) by \(x_h\),
where \(h<\ell\), provided the result remains in the box.
Root-of-unity distractions followed by initial ideals produce an ideal
\(H\supseteq P\) closed under these transfers, with unchanged Hilbert
function and no smaller quotient Betti numbers. In each degree, replace
the box monomials belonging to \(H\) by the same number of greatest
lexicographic monomials in the box, and retain \(P\). Comparing their bounded
shadows---the multiples by one variable that remain in the box---shows
that the resulting space \(G\) is an ideal.

For the Betti comparison, put \(z=x_n\). The modules \(H/zH\) and
\(G/zG\) over the earlier variables have equal Hilbert functions and
retain the Betti numbers of the ideals \(H\) and \(G\). They split
into components according to the exponent of \(z\), from \(0\) to
\(a_n\). Stability makes every earlier variable annihilate the
components at exponents \(1,\ldots,a_n-1\). Only the two extreme
components can therefore contribute nonzero ranks to the Koszul
differentials. The comparisons on the box slices \(u_n=0\) and
\(u_n=a_n-1\), induction in the number of variables, and monotonicity
of differential ranks under submodules and quotients compare these
endpoint contributions. Since the Koszul terms have equal dimensions,
the resulting reverse inequality for adjacent rank sums gives the
required Betti bound.

Sections~\ref{sec:filtered} and \ref{sec:monomial} establish the filtered
comparison and the transfer to \(M\).
Sections~\ref{sec:stabilization}--\ref{sec:tor} give the complete monomial
comparison. Section~\ref{sec:assembly} identifies its output with
\eqref{eq:prescription} and proves Theorem~\ref{thm:main}.
Section~\ref{sec:characteristic-zero} then removes the Artinian and
complex-coefficient restrictions from the Hilbert and Betti conclusions.
